% Part: incompleteness
% Chapter: arithmetization-syntax
% Section: coding-symbols

\documentclass[../../../include/open-logic-section]{subfiles}

\begin{document}

\olfileid{inc}{art}{cod}
\olsection{علامتاں نوں کوڈ دینا}

پہلے درجے دی منطق دی بنیادی زبان~$\Lang L$ ایہ علامتاں ورتدی اے:
\[
\lfalse \quad \lnot \quad \lor \quad \land \quad \lif \quad \lforall
\quad \lexists \quad \eq \quad ( \quad ) \quad ,
\]
نال ایہ متغیراں تے !!{constant}s دے !!{enumerable} مجموعے،
تے دلیلاں دی ہر ممکن تعداد لئی !!{function}s تے !!{predicate}s
دے !!{enumerable} مجموعے وی ورتدی اے۔ اسیں ہر علامت نوں
\emph{کوڈ} وجوں اک عدد ایس طرح دے سکدے آں کہ ہر علامت دا
اپنا وکھرا کوڈ ہووے تے دو وکھریاں علامتاں دا کوڈ اکّو نہ ہووے۔
ایہ ممکن اے، کیوں جے ساریاں علامتاں دا مجموعہ !!{enumerable}
اے، تے ایس مجموعے تے قدرتی عدداں دے مجموعے وچ !!a{bijection}
اے۔ پر اسیں ایہ وی یقینی بنانا چاہندے آں کہ کوڈ توں علامت تے
اوہدے بارے کجھ معلومات (مثلاً !!a{function} دیاں دلیلاں دی
تعداد) قابلِ حساب طریقے نال واپس لبھ سکّیے۔ ایہ کرن دے کئی
طریقے نیں۔ ایتھے اسیں اک طریقہ دیندے آں جیہڑا بنیادی بازگشتی
فنکشن ورتدا اے۔ (یاد رکھو کہ $\tuple{n_0, \dots, n_k}$ اوہ عدد
اے جیہڑا $n_0$, \dots, $n_k$ دے سلسلے دا کوڈ اے۔)

\begin{defn}
جے $s$ زبان~$\Lang L$ دی اک علامت ہووے، تاں
اوہدا \emph{علامتی کوڈ}~$\scode s$ اینج مقرر کردے آں:
\begin{enumerate}
\item جے $s$ منطقی علامتاں وچوں ہووے، تاں
  $\scode s$ دی قدر ایہ جدول دندا اے:
\[
\begin{array}{cccccccccc}
  \lfalse & \lnot & \lor & \land & \lif & \lforall \\
\tuple{0, 0} & \tuple{0, 1} & \tuple{0, 2} & \tuple{0, 3} &
\tuple{0, 4} & \tuple{0, 5} \\
\lexists & \eq & ( & ) & ,\\
\tuple{0, 6} & \tuple{0, 7} &
\tuple{0, 8} & \tuple{0, 9} & \tuple{0, 10}
\end{array}
\]
\item جے $s$، $i$-واں متغیر $\Obj v_i$ ہووے،
  تاں $\scode s = \tuple{1, i}$۔
\item جے $s$، $i$-واں !!{constant}~$\Obj c_i$
  ہووے، تاں $\scode s = \tuple{2, i}$۔
\item جے $s$، $i$-واں $n$-دلیلی
  !!{function}~$\Obj f_i^n$ ہووے، تاں
  $\scode s = \tuple{3, n, i}$۔
\item جے $s$، $i$-واں $n$-دلیلی
  !!{predicate}~$\Obj P_i^n$ ہووے، تاں
  $\scode s = \tuple{4, n, i}$۔
\end{enumerate}
\end{defn}

\begin{prop}
ایہ تعلق بنیادی بازگشتی نیں:
\begin{enumerate}
\item $\fn{Fn}(x, n)$ اُدوں تے صرف اُدوں جدوں $x$،
  $\Obj f^n_i$ دا کوڈ ہووے، کسے~$i$ لئی۔ سو $x$ اک $n$-دلیلی !!{function} دا کوڈ اے۔
\item $\fn{Pred}(x, n)$ اُدوں تے صرف اُدوں جدوں $x$،
  $\Obj P^n_i$ دا کوڈ ہووے، کسے~$i$ لئی، یا $x$، $\eq$ دا کوڈ ہووے تے $n = 2$۔ سو $x$ اک
  $n$-دلیلی !!{predicate} دا کوڈ اے۔
\end{enumerate}
\end{prop}

\begin{defn}
جے $s_0, \dots, s_{n-1}$ علامتاں دا سلسلہ
ہووے، تاں اوہدا \emph{گوڈل نمبر}
$\tuple{\scode{s_0}, \dots, \scode{s_{n-1}}}$ اے۔
\end{defn}

\begin{explain}
دھیان رکھو کہ \emph{کوڈ} تے \emph{گوڈل نمبر} وکھ وکھ
چیزاں نیں۔ مثلاً متغیر~$\Obj v_5$ دا کوڈ~$\scode{\Obj
  v_5} = \tuple{1, 5} = 2^2\cdot 3^6$ اے۔ پر متغیر~$\Obj v_5$
نوں ٹرم وجوں ویکھیے تاں اوہ علامتاں دا اک سلسلہ وی اے، جیہدی
لمبائی~$1$ اے۔ ایس \emph{گوڈل نمبر}~$\Gn{\Obj v_5}$ دی قدر،
جدوں \emph{ٹرم}~$\Obj v_5$ دی گل ہووے، ایہ اے:
$\tuple{\scode{\Obj v_5}} = 2^{\scode{\Obj
    v_5} + 1} = 2^{2^2\cdot 3^6 + 1}$۔
\end{explain}

\begin{ex}
یاد کرو کہ جے $k_0$, \dots, $k_{n-1}$
عدداں دا سلسلہ ہووے، تاں اولی عدداں دیاں قوتاں
والی کوڈکاری وچ سلسلے دا کوڈ
$\tuple{k_0, \dots, k_{n-1}}$ ایہ اے:
\[
2^{k_0+1}\cdot3^{k_1+1}\cdot \dots \cdot p_{n-1}^{k_{n-1}+1},
\]
ایتھے $p_i$، $i$-واں اولی عدد اے
(گنتی $p_0 = 2$ توں شروع ہوندی اے)۔
سو مثلاً فارمولا $\eq[\Obj v_0][\Obj 0]$،
یا ہور کھول کے ${\eq}(\Obj v_0,\Obj c_0)$،
دا گوڈل نمبر ایہ اے:
\[
\tuple{\scode{\eq},\scode{(},\scode{\Obj v_0},\scode{,}, \scode{\Obj
    c_0},\scode{)}}.
\]
ایتھے $\scode{\eq}$، $\tuple{0,7} = 2^{0+1}\cdot
3^{7+1}$ اے، $\scode{\Obj v_0}$،
$\tuple{1,0} = 2^{1+1}\cdot3^{0+1}$ اے،
تے ایداں ہی اگے۔ ایس لئی $\Gn{=(\Obj v_0,\Obj c_0)}$
ایہ بندا اے:
\begin{multline*}
2^{\scode{=} + 1}\cdot 3^{\scode{(}+1}\cdot 5^{\scode{\Obj v_0}+1}
\cdot 7^{\scode{,} + 1} \cdot 11^{\scode{\Obj c_0}+1} \cdot
13^{\scode{)}+1} = \\
2^{2^1\cdot 3^8 + 1}\cdot 3^{2^1\cdot 3^9+1}\cdot 5^{2^2\cdot 3^1+1}
\cdot 7^{2^1\cdot 3^{11} + 1} \cdot 11^{2^3\cdot3^1+1} \cdot
13^{2^1\cdot3^{10}+1} = \\
2^{13\,123}\cdot 3^{39\,367}\cdot 5^{13}\cdot 7^{354\,295}\cdot11^{25}\cdot13^{118\,099}.
\end{multline*}
\end{ex}

\end{document}
